% Vietnamese translation for OI Public Library
% Source-PDF: 动态规划 DYNAMIC PROGRAMMING/动态规划之状态压缩 - 周伟.pdf
\documentclass[11pt]{article}
\usepackage{oipl-vn}

\newcommand{\Oh}{\mathcal{O}}
\newcommand{\lowbit}{\operatorname{lowbit}}
\newcommand{\popcount}{\operatorname{popcount}}

\title{Nén Trạng Thái}
\author{Zhou Wei\\Trường Trung học 101 Trịnh Châu / Đại học Thiên Tân}
\date{}

\begin{document}
\maketitle

\origtitle{Nén trạng thái}
\sourcepath{Dynamic Programming / State Compression - Zhou Wei}

\begin{abstract}
Nhiều bài toán có giá trị ứng dụng thực tế nhưng rất khó có thuật toán đa thức
cho mọi dữ liệu. Tìm kiếm vét cạn luôn đúng nhưng thường quá chậm. Bài viết
giới thiệu tư tưởng nén trạng thái qua một loạt ví dụ: dùng bit để biểu diễn
tập hợp hoặc cấu hình cục bộ, rồi xây dựng truy hồi hoặc quy hoạch động trên
những trạng thái đã nén.
\end{abstract}

\paragraph{Từ khóa.}
Nén trạng thái; băm; quy hoạch động; truy hồi; thao tác bit.

\tableofcontents

\section{Mở đầu}

Trong OI, ta gặp nhiều thuật toán đa thức như tìm kiếm nhị phân, Euclid, cây
cân bằng, đường đi ngắn nhất. Tuy vậy cũng có những lớp bài như NPC hoặc
NP-hard mà chưa biết thuật toán đa thức tổng quát: Hamilton cycle, TSP, đường
đi đơn dài nhất trong đồ thị có hướng, v.v. Với những bài như vậy, nếu chỉ
tìm kiếm thuần túy thì độ phức tạp có thể là \(n!\) hoặc tệ hơn. Nén trạng
thái là một cách thường dùng để có thuật toán mũ nhưng nhỏ hơn và dễ cài đặt
hơn nhiều so với vét cạn trực tiếp.

\subsection{Thao tác bit cơ bản}

Các phép bit thường gặp:
\[
\begin{array}{c|c|c|c}
\text{Tên} & \text{C/C++} & \text{Pascal} & \text{Ý nghĩa}\\ \hline
\text{AND} & \& & \texttt{and} & \text{tất cả là }1\text{ thì }1\\
\text{OR} & | & \texttt{or} & \text{có }1\text{ thì }1\\
\text{NOT} & \sim & \texttt{not} & 0\leftrightarrow 1\\
\text{XOR} & \^{} & \texttt{xor} & \text{khác nhau thì }1
\end{array}
\]
Một số ứng dụng nhỏ:
\begin{itemize}
  \item Lấy một số bit của \(x\): dùng AND với mặt nạ.
  \item Lấy bit \(1\) thấp nhất:
  \[
    \lowbit(x)=x\mathbin{\&}(-x).
  \]
  \item Đặt một số bit thành \(1\): dùng OR với mặt nạ.
  \item Đảo một số bit: dùng XOR với mặt nạ.
\end{itemize}

\section{Khởi động: đặt xe trên bàn cờ}

\paragraph{Bài dẫn.}
Trên bàn cờ \(n\times n\), \(n\le 20\), đặt \(n\) quân xe sao cho không quân
nào ăn quân nào. Đếm số cách đặt.

Đây là một bài tổ hợp đơn giản, đáp án là \(n!\). Nhưng để giới thiệu nén
trạng thái, ta giải theo hướng truy hồi. Đặt từng hàng một; một trạng thái
là tập các cột đã có quân xe. Nếu cột đã dùng thì bit tương ứng là \(1\),
ngược lại là \(0\). Ví dụ với \(n=5\), trạng thái \(01101_2\) nghĩa là các
cột \(1,3,4\) đã dùng.

Gọi \(f_s\) là số cách đạt tới trạng thái \(s\). Khi \(s=01101_2\), trạng
thái này có thể đến từ ba trạng thái trước đó bằng cách bỏ lần lượt bit ở
cột \(1,3,4\):
\[
  f_{01101}=f_{01100}+f_{01001}+f_{00101}.
\]
Tổng quát:
\[
  f_0=1,\qquad
  f_s=\sum_{i:\text{ bit }i\text{ của }s=1} f_{s-2^i}.
\]
Thuật toán có độ phức tạp \(\Oh(n2^n)\) và bộ nhớ \(\Oh(2^n)\). Nó kém hơn
công thức \(n!\), nhưng có khả năng mở rộng tốt hơn.

\section{Mô hình bàn cờ}

\subsection{Xe với ô cấm}

Nếu một số ô không được đặt quân xe, công thức tổ hợp không còn rõ ràng, nhưng
truy hồi nén trạng thái chỉ cần sửa rất ít. Với hàng \(r\), dùng mặt nạ
\(a_r\) để đánh dấu các cột có thể đặt. Khi trạng thái hiện tại là \(s\), ta
chỉ liệt kê các bit \(1\) trong
\[
  s\mathbin{\&} a_r,
\]
nhờ vậy không cần kiểm tra từng cột trong vòng lặp. Đây là ví dụ điển hình
cho việc dùng bit mask để biến ràng buộc cục bộ thành phép toán hằng số.

\subsection{Đặt quân không kề nhau}

Cho bàn cờ \(n\times m\), \(n,m\le 80\), \(nm\le 80\). Đặt \(k\le 20\) quân
sao cho không có hai quân kề nhau theo hàng hoặc cột. Vì \(nm\le 80\), ít
nhất một trong hai chiều không vượt quá \(8\); chọn chiều đó làm số cột.

Ta tiền xử lý mọi trạng thái hợp lệ của một hàng vào mảng \(s_1,\ldots,s_t\),
tức là các mask không có hai bit \(1\) liền nhau. Gọi \(c_j\) là số bit \(1\)
của trạng thái \(s_j\). Đặt
\[
  f_{i,j,k}
\]
là số cách xét xong \(i\) hàng, hàng \(i\) dùng trạng thái \(s_j\), và đã đặt
tổng cộng \(k\) quân. Khi hàng hiện tại \(s_j\) không xung đột với hàng trước
\(s_p\), tức
\[
  s_j\mathbin{\&}s_p=0,
\]
ta có
\[
  f_{i,j,k}=\sum_p f_{i-1,p,k-c_j}.
\]
Có thể dùng mảng cuộn vì chỉ phụ thuộc hàng trước.

\subsection{Little Kings}

Với bài đặt \(k\) quân vua trên bàn \(n\times n\), \(n\le 10\), quân vua ăn
theo cả tám hướng. Trạng thái một hàng vẫn là mask không có hai bit \(1\) kề
nhau. Khác biệt là hàng hiện tại còn không được chéo với hàng trước. Có thể
tiền xử lý một mask \(q_j\) biểu diễn những ô ở hàng trước bị cấm khi hàng
hiện tại là \(s_j\). Điều kiện tương thích là
\[
  q_j\mathbin{\&}s_p=0.
\]
Công thức DP giữ nguyên dạng như trên.

\subsection{Pháo binh}

Trong bài \emph{Artillery Positions}, bàn cờ \(n\times m\), \(n\le100\),
\(m\le10\), có một số ô cấm. Cần đặt nhiều quân nhất sao cho trên cùng hàng
hoặc cột, hai quân cách nhau ít nhất hai ô trống.

Tầm ảnh hưởng theo hàng/cột là \(2\), nên trạng thái cần nhớ hai hàng. Gọi
\[
  f_{i,j,k}
\]
là số quân tối đa sau \(i\) hàng, trong đó hàng \(i\) có trạng thái \(s_j\)
và hàng \(i-1\) có trạng thái \(s_k\). Khi ba trạng thái \(s_j,s_k,s_l\) đôi
một không xung đột và \(s_j\) không đặt vào ô cấm, ta chuyển:
\[
  f_{i,j,k}=\max_l f_{i-1,k,l}+c_j.
\]
Vì số trạng thái hợp lệ của một hàng nhỏ, thuật toán \(\Oh(n t^3)\) đủ nhanh.
Có thể tiền xử lý các cặp hoặc bộ trạng thái tương thích để giảm hằng số.

\section{Mô hình phủ}

\subsection{Domino 1 x 2}

Cho bàn \(n\times m\), \(1\le n,m\le11\), dùng domino \(1\times2\) phủ kín
bàn. Giả sử \(m\le n\). Gọi \(f_{i,s}\) là số cách đã phủ kín \(i-1\) hàng,
và hàng \(i\) đang có trạng thái phủ \(s\). Ta DFS mọi cách đặt domino ở hàng
hiện tại để sinh cặp trạng thái \((s_1,s_2)\): trạng thái để lại cho hàng
hiện tại và trạng thái yêu cầu từ hàng trước. Công thức:
\[
  f_{0,2^m-1}=1,\qquad
  f_{i,s_1}=\sum f_{i-1,s_2},
\]
trong đó \((s_1,s_2)\) là một cách đặt hợp lệ.

Cách DFS tự nhiên có ba lựa chọn tại một vị trí: đặt domino dọc, đặt domino
ngang, hoặc không đặt tại ô đó tùy theo trạng thái đã phủ. Số cách sinh ra
không phải \(4^m\); nó thỏa truy hồi gần với
\[
  h_0=1,\qquad h_1=2,\qquad h_i=2h_{i-1}+h_{i-2},
\]
nên thực tế nhỏ hơn rất nhiều so với liệt kê thô.

\subsection{Tăng tốc bằng lũy thừa ma trận}

Nếu \(n\) rất lớn, DP theo từng hàng sẽ chậm. Ta xây đồ thị có \(2^m\) đỉnh,
mỗi đỉnh là một trạng thái hàng. Nếu \((s_1,s_2)\) là một chuyển hợp lệ, thêm
cạnh từ \(s_2\) tới \(s_1\). Khi đó số cách đi từ trạng thái đầy đủ
\((2^m-1)\) về chính nó sau \(n\) bước là đáp án. Với ma trận kề \(G\), đáp án
là phần tử tương ứng của \(G^n\).

Ma trận này rất thưa, nên có thể lưu các phần tử khác \(0\) bằng bộ ba
\((p,q,value)\) và nhân ma trận thưa: với mỗi phần tử khác \(0\) của \(A\),
tìm các phần tử cùng hàng phù hợp trong \(B\), rồi cộng vào kết quả. Nhờ đó
trong nhiều trường hợp thực tế tốt hơn rất nhiều so với nhân ma trận đặc
\(\Oh(8^m)\).

\subsection{Các quân phủ khác}

Với domino \(1\times2\) và quân chữ L \(2\times2\) bỏ một ô, trạng thái và
công thức giống domino, nhưng DFS sinh chuyển có nhiều trường hợp hơn. Với
quân \(1\times r\), trạng thái có thể phải ghi nhớ \(r-1\) hàng. Nếu lưu trực
tiếp nhiều mask nhị phân sẽ có nhiều trạng thái dư thừa; tốt hơn là dùng hệ
cơ số \(r\). Mỗi chữ số mô tả số ô liên tiếp đã được phủ theo chiều dọc tại
cột tương ứng. Đây là ví dụ cho thấy ``nén trạng thái'' không nhất thiết chỉ
là nhị phân.

\subsection{Bugs Integrated}

Với bàn \(n\times m\), \(n\le150\), \(m\le10\), cần đặt nhiều quân \(2\times3\)
nhất, một số ô bị cấm. Dùng hệ tam phân để biểu diễn trạng thái của hai hàng:
chữ số \(0,1,2\) mô tả mức đã phủ liên tục từ hàng trên xuống tại cột đó.
Đặt \(f_{i,s}\) là số quân tối đa tới hàng \(i\) với trạng thái \(s\). DFS tại
từng hàng sinh chuyển từ trạng thái hàng trước sang hàng hiện tại, thử đặt
quân ngang hoặc dọc nếu không chạm ô cấm, đồng thời có thể đặt các ``ô ảo''
để đóng những vị trí không dùng. Công thức tổng quát:
\[
  f_{i,s_1}=\max_s \{f_{i-1,s}+num(s\to s_1)\}.
\]

\section{Bản chất của nén trạng thái}

Các ví dụ trên đều là những truy hồi hoặc DP thông thường, chỉ có một hoặc
nhiều chiều trạng thái được ``nén'' thành số nguyên. Có thể xem đây là một
dạng băm: một tập hợp, một cấu hình hàng, hoặc một trạng thái phủ được mã hóa
bằng bit hoặc chữ số trong một cơ số nhỏ.

Động cơ của nén trạng thái thường là: trạng thái quá thô không đủ để chuyển
đúng. Ví dụ, chỉ biết đã xử lý tới hàng \(i\) là không đủ để phủ domino, vì
hàng kế tiếp còn bị ảnh hưởng bởi các ô nhô xuống từ hàng trước. Ta cần ghi
thêm cấu hình biên; nén trạng thái giúp cấu hình đó đủ nhỏ để liệt kê.

\section{Ví dụ đồ thị}

\subsection{Đếm ghép đầy đủ trong đồ thị hai phía}

Có \(n\le20\) con bò và \(m\le20\) đồng cỏ. Mỗi con bò thích một số đồng cỏ
và không muốn chung đồng cỏ với bò khác. Hỏi có bao nhiêu cách gán mỗi con bò
một đồng cỏ nó thích.

Dùng một mask \(s\) biểu diễn tập đồng cỏ đã dùng. Gọi \(z(s)=\popcount(s)\)
là số bò đã xử lý. Nếu ma trận kề là \(g\), trong đó \(g_{i,j}=1\) khi bò
\(i\) thích đồng cỏ \(j\), thì
\[
  f_0=1,\qquad
  f_s=\sum_{j:\text{ bit }j\text{ của }s=1}
       f_{s-2^j}\,g_{z(s),j}.
\]
Đáp án là tổng \(f_s\) trên các trạng thái có \(\popcount(s)=n\). Có thể hiểu
tương đương theo tập hợp: \(f_S\) là số cách gán các bò đầu tiên vào đúng các
đồng cỏ thuộc tập \(S\); để thu nhỏ bài toán, liệt kê đồng cỏ của con bò cuối.

\subsection{Đếm số thứ tự topo}

Cho đồ thị có hướng \(n\le20\), cần đếm số thứ tự topo hợp lệ. Trạng thái
\(S\) là tập đỉnh đã đặt vào thứ tự. \(f_S\) là số thứ tự topo của các đỉnh
trong \(S\). Ta có thể liệt kê đỉnh cuối cùng của \(S\); đỉnh đó hợp lệ nếu
mọi cạnh đi ra từ nó tới trong \(S\) đều đã được xử lý phù hợp, hoặc tương
đương khi bỏ nó khỏi \(S\), các điều kiện tiền nhiệm vẫn đúng. Chi tiết phụ
thuộc cách lưu cạnh, nhưng mô hình vẫn là DP tập con kinh điển.

\subsection{Shortest Common Superstring}

Cho \(n\le15\) xâu, cần tìm xâu ngắn nhất chứa mọi xâu đã cho làm xâu con;
nếu có nhiều đáp án, chọn từ điển nhỏ nhất. Trước hết loại xâu nào là xâu con
của xâu khác. Sau đó tiền xử lý độ chồng lớn nhất giữa mọi cặp xâu. Trạng thái
\[
  f_{S,i}
\]
là xâu ngắn nhất tạo từ tập \(S\) và kết thúc bằng xâu \(i\). Chuyển từ
\(f_{S-\{i\},j}\) sang \(f_{S,i}\) bằng cách nối thêm phần chưa chồng của
\(i\). Đây là một ví dụ khác: trạng thái là tập xâu đã dùng, nén bằng mask.

\section{Kết luận}

Nén trạng thái không phải là một thuật toán riêng lẻ, mà là cách biểu diễn
trạng thái. Khi số phần tử nhỏ, hoặc khi chỉ cần nhớ một ``biên'' hẹp của bài
toán, bit mask và cơ số nhỏ giúp biến tìm kiếm thô thành truy hồi hoặc quy
hoạch động có cấu trúc. Điểm quan trọng là chọn đúng phần thông tin cần lưu:
đủ để chuyển trạng thái không có hậu hiệu, nhưng không dư thừa đến mức làm
nổ số trạng thái.

\end{document}
